\documentclass[11pt]{article}
\usepackage[margin=1in]{geometry}
\usepackage{amsmath,amssymb,amsthm,array,longtable,hyperref}
\newtheorem{theorem}{Theorem}
\newtheorem{lemma}{Lemma}
\newcommand{\Z}{\mathbb Z}
\newcommand{\ind}{\mathbf1}
\title{Doubled staircases and their near-one rotations:\\replacement base recurrence proofs}
\author{Audited integration of the 20 October base proofs}
\date{26 October 2026}
\begin{document}\maketitle

\begin{sloppypar}
\paragraph{Scope and provenance.}
The first two theorems replace Lemmas 6.15 and 6.16 of
\texttt{revised\_version2/farey\_recursion.tex}. Their explicit ranks and
uniform case proofs were already established in
\path{flattice-thm13-20261019/BASE-B5-ANALYSIS.md}, section 7, and
\texttt{BASE-B7-RANK.md}. The latter has the separate
\texttt{BASE-B7-RANK-AUDIT.md}. The period extensions and complete
Laplacian tables used below are prior results of
\texttt{BASE-B5-ANALYSIS.md}, sections 1--5, and
\texttt{BASE-B7-ANALYSIS.md}; the harmonic rotation is from
\texttt{BASE-DUALITY-PROOF.md}. These files lie under
\path{thm13-recurrence-agents/flattice-thm13-20261019/}.
This is an integration of proved base results, not a new discovery or an
arbitrary-word recurrence argument. The exact imported extension facts
are stated below; the decreasing-rank proofs and phase transfer are
included. The final section corrects a domain typo and a necessary phase
convention in the rotated alternate statement, Lemma 6.10.
\end{sloppypar}

Throughout $d\ge4$ is even. Set
\[
 T=d^2+2d-1,\quad U=(2d,2),\quad V=(1-d,d+1),\quad L=\Z U+\Z V,
 \quad p=(x+y)\bmod2,\quad t=x+y.
\]
Bit $0$ denotes the original graph: inputs are horizontal at even
coordinate sum and vertical at odd sum. Bit $1$ reverses these choices.
For $s=1+\Delta g\in\{0,1\}$, a forbidden set must contain both inputs
at an occupied $s=1$ point and at least one input at an occupied $s=0$
point. Consequently an integer rank decreasing to one input when $s=1$
and to both when $s=0$ excludes every nonempty finite forbidden set.
This is the forbidden-set characterization of recurrence, not a
uniqueness or maximality hypothesis.

\section*{1. Standard doubled staircase: the exact base data}
Let $D_5$ be the integer points in
\[
 2-d\le x\le2d,\quad0\le y\le d+2,\quad1\le x+y\le2d+3,
\]
with $y=0,x\ge d$ and $y=d+2,x\le2$ deleted, and with
$a=(0,0)$ and $b=(d+2,d+2)$ added. Put
\[
 J_5=\ind_{\{x+y=d+2,\ 2\le y\le d\}\cup\{a,b\}},
\]\[
 g_5=-\frac{y(y-1)}2+\min(x-1,0)+\min(d+1-x,0)
            +\max(x+y-d-2,0)+J_5.                         \tag{1}
\]
These are exactly the paper's standard tile and values. The prior
extension calculation proves that $D_5$ represents every $L$ class and
that (1) has a unique compatible integer extension with
\[
 g_5(z+U)-g_5(z)=0,\qquad
 g_5(z+V)-g_5(z)=x-dy-\frac{d(d+3)}2.                       \tag{2}
\]
Its growth is $T^{-1}\left(\begin{smallmatrix}-1&d\\d&-d^2\end{smallmatrix}\right)$.
For precision, the exhaustive positive overlap sources are
\[
\begin{array}{c|l}
U&a;\ (1-y,y),\ 1\le y\le d-1;\ (2-d,d)\\
V&(x,0),\ 1\le x\le d-1;\ (x,1),\ d+2\le x\le2d.
\end{array}                                                   \tag{3}
\]
There are no other nonzero overlaps, up to reversing endpoints, and no
triple classes. The tile has $2d^2+7d-3$ points; subtracting its
$3d-1$ repeated pairs gives $2T=\det(U,V)$ classes. Thus (3), together
with the directly substituted increments (2), also explains the covering
and compatibility dependency without a hidden tiling assertion.

The extension is binary superharmonic. Its even $s=0$ points in $D_5$
are exactly
\[
\begin{split}
 A_1&=\{(1,y):1\le y\le d-1,\ y\text{ odd}\},\\
 A_2&=\{(d+1,y):3\le y\le d+1,\ y\text{ odd}\},\\
 A_3&=\{(x,y):x+y=d+2,\ 2\le y\le d\}.
\end{split}                                                   \tag{4}
\]
To state its odd $s=1$ points, set $w=(d+2,d+2)$ and
\[
\begin{split}
 B_-={}&\{(x,1-x):2-d\le x\le-1\}
 \cup\{(x,0):1\le x\le d-1,\ x\text{ odd}\}\\
 &\cup\{(x,1):d\le x\le2d,\ x\text{ even}\},\qquad B_+=w-B_-,\\
 Z_1={}&\{x+y=d+1,\ 1\le y\le d-1\},\quad
 Z_2=\{x+y=d+3,\ 3\le y\le d+1\},\\
 Z_3={}&\{(0,1),(d+2,d+1)\}.
\end{split}
\]
The odd $s=1$ set is exactly $B_-\cup B_+\cup Z_1\cup Z_2\cup Z_3$;
all remaining odd points have $s=0$. These complete tables include
physical stars crossing the tile boundary. They are obtained in the
cited base analysis by direct second differences of (1), substituting
(2) at exits. In particular the subsequent proof imports no unproved
statement about recurrence on embedded copies in a different field.

\section*{2. Standard decreasing rank}
Write $c=(2d,2)$, $e=(2d,1)$, $f=(d+1,d+2)$. Define $\rho_5$ on $D_5$
by the first applicable clause:
\[
\begin{array}{c|l}
\rho_5&\text{condition}\\\hline
4&z=a\text{ or }c\\
5&z=e\text{ or }f\\
0&z=b\\
p&y=0\text{ or }x\le0\\
8+p&y=d+2,\ \text{or }y=1,x\ge d+2,\ \text{or }x\le d,t\ge d+3\\
8(d+1-x)+p&t\le d+1\\
8&t=d+2,\ x\in\{1,d+1\}\\
8d&t=d+2,\ x\notin\{1,d+1\}\\
8(2d-x)+p&t\le2d+1\\
8&t=2d+2\\
1&t=2d+3.
\end{array}                                                    \tag{5}
\]
The clauses cover $D_5$. On the successive $U$ overlaps in (3), the
paired values are $4$, $1$, and $0$. On the bottom $V$ overlaps they
are $p$; on the row-one $V$ overlaps they are $8+p$, except for the last
pair $e,f$, where they are $5$. Therefore $\rho_5$ extends periodically.

\begin{theorem}[Replacement for Lemma 6.15]
For every even $d\ge4$, the extension (1)--(2) is integer superharmonic
and recurrent. The rank $R_5(x,y)=6x+\rho_5(x,y)$ decreases by at least
one to both inputs at $s=0$ and to one specified input at $s=1$.
\end{theorem}
\begin{proof}
For even $s=0$ points, the two differences
$R_5(\text{input})-R_5(\text{center})$, in left--right order, are
\[
\begin{array}{c|c}
A_1&-6+1-8d,\ -1\\
A_2,\ y\le d-1&3-8(d-1),\ -1\\
(d+1,d+1)&-5,\ -1\\
A_3,\ x\le d-1&11-8x,\ 15-8d\\
A_3,\ x=d&11-8d,\ -1.
\end{array}                                                    \tag{6}
\]
All are negative, and all these inputs lie inside the tile.
For even $s=1$ points use the following exhaustive ordered list, after
removing (4). Each chosen move has difference at most $-1$.
\[
\begin{array}{l|c|c}
\text{center}&\text{move}&\text{difference}\\\hline
 a,c,b,(2-d,d)&\text{left}&-1\\
 x\le0\text{ otherwise}&\text{left}&-5\\
 y=0,\ 2\le x\le d-2&\text{left}&-5\\
 y=d+2,\ 4\le x\le d&\text{left}&-5\\
 y=1,\ d+3\le x\le2d-1&\text{left}&-5\\
 x\ge2,\ y\ge1,\ t\le d&\text{right}&-1\\
 (1,d+1)&\text{left}&-13\\
 (d+1,1)&\text{left}&-5\\
 x\le d,\ y\le d+1,\ t\ge d+4&\text{left}&-5\\
 x\ge d+2,\ y\ge2,\ t\le2d&\text{right}&-1\\
 t=2d+2,\ x\ge d+2,\ y\ge3&\text{right}&-1.
\end{array}                                                    \tag{7}
\]
The only outside chosen inputs are the left input of $a$, represented
by $(2d-1,2)$ with correction $9$, and the left input of $(2-d,d)$,
represented by $f$ with correction $5$. Actual input coordinates,
not representative coordinates, are used in the linear term $6x$.
Exhaustiveness follows by first removing the exceptional points, the
nonpositive region and the bottom/top/shelf rows. The remaining even
sums are $t\le d$, $t=d+2$, and $t\ge d+4$; the last region splits at
$x=d$, with spine $x=d+1$ already in (4).

Every odd $s=0$ point has both inputs inside $D_5$. There are four
possibilities. If $x\le0$, its correction is $1$ and both inputs have
correction zero; the only exceptional possible input $a$ would require
center $(0,1)$, which has $s=1$. In the left high region, the center
correction is $8(d+1-x)+1$ and both inputs have correction at most
$8(d+1-x)$. The only $s=0$ center on $t=d+1$ is $(1,d)$; its upper
input has correction $8\le8d$. In the middle region, an odd $s=0$
center has $t\ge d+5$, correction $9$, and both input corrections $8$.
In the right high region its correction is $8(2d-x)+1$ and both inputs
have correction at most $8(2d-x)$; the possible shelf or diagonal input
has correction $8$, with $x\le2d-1$. The exceptional point $c$ cannot
be such an input. Thus both vertical rank differences are negative.

For odd $s=1$ points the following table supplies a correction exactly
one smaller. Entries are read in order. The two repetitions
$B_-\cap Z_1=\{(d,1)\}$ and $B_+\cap Z_2=\{(2,d+1)\}$ are omitted
from the later diagonal rows.
\begin{center}\small
\begin{longtable}{lcl}
center&move&outside representative, when required\\\hline
$(x,1-x),\ 2-d\le x\le-1$&up&\\
$(x,0),\ x\text{ odd},\ 1\le x\le d-1$&down&$(x+1-d,d)$\\
$(d,1)$&down&$(1,d+1)$\\
$(x,1),\ x\text{ even},\ d+2\le x\le2d-2$&down&$(x+1-d,d+1)$\\
$e$&up&\\
$t=2d+3,\ d+3\le x\le2d$&up&$(x-2d,y-1)$\\
$(x,d+2),\ x\text{ odd},\ 3\le x\le d-1$&down&\\
$f$&up&$c$\\
$(x,d+1),\ x\text{ even},\ 2-d\le x\le0$&down&\\
$(2,d+1)$&up&$(d+1,1)$\\
$Z_1\setminus\{(d,1)\}$&down&\\
$Z_2\setminus\{(2,d+1)\}$&up&\\
$(0,1),(d+2,d+1)$&up&
\end{longtable}\end{center}
The correction differences follow immediately from (5). These are all
odd $s=1$ points by the displayed exact classification. Periodicity
transports every physical star inequality to the plane. A minimal-rank
point of a finite nonempty forbidden set now gives a contradiction.
\end{proof}

\section*{3. Alternate doubled staircase: the exact base data}
Let $D_7$ be the integer points in
\[
 1-d\le x\le2d-2,\quad-1\le y\le d+2,\quad-1\le x+y\le2d+1,
\]
with $y\le0,x\ge d$ and $y\ge d+1,x\le-1$ deleted, and with
$a=(-1,-1)$, $b=(d,d+2)$ added. Put
\[
 J_7=\ind_{\{x+y=d,\ 1\le x\le d-2\}\cup\{a,b\}},
\]\[
 g_7=-\frac{y(y+1)}2+\min(x,0)+\min(d-1-x,0)
                  +\max(x+y-d,0)+J_7.                     \tag{8}
\]
These agree literally with the paper's alternate tile and heights.
The prior complete extension proof gives
\[
 g_7(z+U)-g_7(z)=-1,\qquad
 g_7(z+V)-g_7(z)=x-dy-\frac{d^2+5d-2}{2},                 \tag{9}
\]
and growth $M(1,d)$. The exhaustive positive overlap sources are
\[
\begin{array}{c|l}
U&(-1-y,y),\ 1\le y\le d-2\\
V&(j,-1),\ 0\le j\le d-1;\ (d-1,0);\ (j,1),\ d-1\le j\le2d-2\\
U+V&a.
\end{array}                                                   \tag{10}
\]
Only these vectors and their negatives occur. There are no triple classes:
a possible triangle would contain the unique $U+V$ pair $a,b$, but both
candidate third points are outside the tile. The tile has $2d^2+7d-2$
points and $3d$ repeated pairs, hence represents all $2T$ classes.
Substitution in (8) gives (9) on all overlaps. These facts justify the
compatible extension independently of any recurrence assertion.

The exact binary Laplacian table has the even $s=0$ set
\[
 \{x=0,\ 0\le y\le d-2\}\ \cup\
 \{x=d-1,\ 3\le y\le d+1\}\ \cup\
 \{x+y=d,\ 1\le x\le d-2\},                              \tag{11}
\]
with even parity imposed. The odd $s=1$ set comprises the following
disjoint lists, with odd parity imposed:
\[
\begin{array}{ll}
y=-1,\ 0\le x\le d-2;&y=d+2,\ 1\le x\le d-1;\\
x+y=-1,\ 1\le y\le d-2;&x+y=2d+1,\ d+1\le x\le2d-2;\\
y=1,\ d\le x\le2d-2;&y=d,\ 1-d\le x\le-1;\\
x+y=d-1\text{ or }d+1,\ 1\le x\le d-2;&(-1,0),(d,d+1).
\end{array}                                                   \tag{12}
\]
All remaining even points have $s=1$, and remaining odd points $s=0$.
The imported proof of (11)--(12) uses the ordinary second differences
of (8), its even correction points and the full boundary-exit table,
with (9) supplying every outside value. It covers $d=4$, where the
middle range of the row table is empty.

\section*{4. Alternate decreasing rank}
Define $\rho_7$ by these disjoint regions in $D_7$:
\[
\begin{array}{c|l|l}
\text{region}&\text{condition}&\rho_7\\\hline
C&z=a\text{ or }b&4\\
E&z=(d,d+1)&5\\
T_0&y=d+2,\ 0\le x\le d-1&8(d-x)+p\\
S&y=d+1,\ 0\le x\le d-1&8d-4x+p\\
B&x+y=2d,\ d\le x\le2d-2&8\\
D_0&x+y=d,\ 1\le x\le d-2&8(d+y)-4\\
L_0&x\ge0,\ y\ge0,\ x+y\le d-1&8(2d-1-x)+p\\
R_0&x\ge d-1,\ 1\le y\le d,\ x+y\le2d-1&8(2d-1-x)+p\\
M&1\le x\le d-2,\ y\le d,\ x+y\ge d+1&8d-4x+p\\
Z&\text{all other points}&p.
\end{array}                                                    \tag{13}
\]
On the $U$ overlaps both corrections are $1$. The bottom-row $V$
overlaps join two $Z$ points, with correction $p$; the isolated $V$
pair joins $(d-1,0)$ to $(0,d+1)$, both of correction $8d+1$.
The row-one $V$ pairs join $R_0$ to $T_0$, with common correction
$8(2d-1-j)+p$. The $U+V$ pair has correction $4$.
Thus (10) gives a well-defined periodic extension of $\rho_7$.

\begin{theorem}[Replacement for Lemma 6.16]
For every even $d\ge4$, the extension (8)--(9) is integer superharmonic
and recurrent. The rank $R_7(x,y)=6x+\rho_7(x,y)$ has the same one/both
decreasing-input property as $R_5$.
\end{theorem}
\begin{proof}
At the even $s=0$ points of (11), the left--right rank differences are
\[
\begin{array}{c|c}
x=0&3-16d,\ -1\\
x=d-1,\ y\le d-1&3-4d,\ -1\\
(d-1,d+1)&-1,\ 7-4d\\
x+y=d,\ x\le d-3&-1,\ 4x-8d+7\\
(d-2,2)&-1,\ -5.
\end{array}                                                    \tag{14}
\]
They are strictly negative even at the minimum parameter $d=4$.
For even $s=1$ points choose left in $C,S,M,Z$ and right in
$T_0,B,L_0,R_0$. There are no such points in $D_0$, and $E$ is odd.
In each high region the difference is $-1$; in $Z$ it is $-5$, except
at the sole left exit $(1-d,d-1)$, where it is $-1$.
The complete outside choices are: from $a$, its left input is represented
by $(2d-2,1)$ with correction $9$; from $(2d-2,2)$ in $B$, its right
input is represented by $(-1,0)$ with correction $1$; and the just
mentioned $Z$ exit is represented by $E$, with correction $5$.
At $b$ the left correction is $9$. The region inequalities ensure that
all other chosen inputs stay in their stated region. For example even
$M$ points satisfy $x\ge2$, so their left input remains in $M$.

Odd $s=0$ points lie only in $S,L_0,R_0,M,Z$. In $S$, their correction
is $8d-4x+1$. At $x=0$ the lower and upper corrections are $0,8d$;
at $x\ge2$ they are $8d-4x,8d-8x$, both smaller. In $L_0$ put
$H=16d-8-8x$. Every $s=0$ center has correction $H+1$ and its inputs
are $L_0$ of correction $H$, $Z$ of correction zero, or $(d-1,1)$ of
correction $H=8d$. A possible transition to $D_0$ would lie on the
excluded $s=1$ diagonal in (12).
In $R_0$ the same $H$ applies. An odd $s=0$ point has $y\ge2$; the
lower input has correction $H$. Its upper input has correction $H$
unless it crosses to $B$, of correction $8\le H$, or to $(d-1,d+1)$
in $S$, of correction $4d+4\le8d=H$. The former inequality uses
$x\le2d-3$, forced by $y\ge2$ and odd sum at most $2d-1$.
In $M$, excluding the $s=1$ diagonal gives $x+y\ge d+3$, so both inputs
are in $M$ or $S$ with correction $8d-4x$, one less than the center.
Finally, an odd $Z$ point with $s=0$ has $x<0$, $1\le y\le d-1$;
(12) excludes every exit and exceptional neighbor. Both inputs have
correction zero, one less than its correction $1$.

The odd $s=1$ cases are exactly these choices; every chosen correction
is one less, and every move is vertical:
\begin{center}\small
\begin{longtable}{lcl}
center&move&outside representative, when required\\\hline
$E=(d,d+1)$&up&\\
$T_0:\ y=d+2,\ x\text{ odd},\ 1\le x\le d-1$&up&$(x+d-1,2)$\\
$L_0:\ x+y=d-1,\ 1\le x\le d-2$&down&\\
$R_0:\ y=1,\ x\text{ even},\ d\le x\le2d-2$&up&\\
$M:\ x+y=d+1,\ 1\le x\le d-2$&up&\\
$Z:\ y=-1,\ x\text{ even},\ 0\le x\le d-2$&down&$(x+1-d,d-1)$\\
$Z:\ x+y=-1,\ 1\le y\le d-2$&up&\\
$Z:\ x+y=2d+1,\ d+1\le x\le2d-2$&up&$(x-2d,y-1)$\\
$Z:\ y=d,\ x\text{ odd},\ 1-d\le x\le-1$&down&\\
$Z:\ (-1,0)$&up&
\end{longtable}\end{center}
These statements follow by substitution in (13), including all endpoints.
For example the last odd $T_0$ point maps to a $B$ representative of
correction $8$, while the others map to $R_0$; and the first $M$
diagonal point moves into $S$ with the same stated correction.
Periodicity extends all inequalities globally. The minimal-rank
forbidden-set argument proves recurrence.
\end{proof}

\section*{5. Stronger directional consequences and endpoints}
The corrections satisfy
\[
 0\le\rho_5\le8d+1,\qquad 0\le\rho_7\le16d-7.
\]
Thus neither field admits a nonempty forbidden set whose $x$ coordinate
is bounded below, even if the set is infinite. For a domain of finite
horizontal span $W$, synchronous burning from its complement finishes
within $6W+8d+2$ rounds for B5, and $6W+16d-6$ for B7. Indeed the
integer rank range has that length minus one, and a point surviving
$t$ rounds forces a decreasing sequence of length $t$ inside the domain.
This also applies to an infinite vertical strip.

No formula here asserts B5 or B7 at $d=2$. In the actual algorithm the
$1/2$ outputs belong to the separately defined B2/B4 families (duals of
the $d=3$ ordinary staircases); the near-one value at $d=2$ is $1/3$,
already covered by the ordinary staircase theorem. The minimum doubled
parameter is $d=4$; every interval and inequality above is valid there.
The correspondence to pure words is $d=2k+2$ with $k\ge1$.

\section*{6. Harmonic rotation, with the physical phase retained}
Set
\[
 \mathcal R(x,y)=(y,-x),\qquad
 K(X,Y)=\frac{X(X-1)}2-\left\lfloor\frac{(X-Y)^2}{4}\right\rfloor,
 \qquad H=\frac12\begin{pmatrix}1&1\\1&-1\end{pmatrix}.
\]
Writing $p=(X+Y)\bmod2$ gives
$4K=X^2+2XY-Y^2-2X+p$. At an odd center the horizontal second
difference is $1/2-1/2=0$; at an even center the vertical one is
$-1/2+1/2=0$. Hence $K$ is harmonic on bit $1$, with growth matrix $H$.
An even translation after $\mathcal R$ preserves coordinate parity,
whereas the rotation interchanges horizontal and vertical directions.
It is therefore a directed graph isomorphism from bit $0$ to bit $1$,
not from bit $0$ to itself.

Define, in formula coordinates,
\[
\begin{split}
 F_6(X,Y)&=g_5(2d+1-Y,X+1)+K(X,Y),\\
 F_8(X,Y)&=g_7(2d-1-Y,X+1)+K(X,Y)+X+1.
\end{split}                                              \tag{15}
\]
The translation vectors in these rotations are $(-1,2d+1)$ and
$(-1,2d-1)$, both even. Thus $\Delta_1F_i$ is the corresponding
$\Delta_0g_j$ composed with the inverse rotation. Affine additions and
$K$ do not change this identity, so the complete forbidden-set rule,
including both incoming alternatives, is transported bijectively.
Consequently $F_6,F_8$ are recurrent superharmonic functions on bit $1$.
The ranks are exactly $R_j$ composed with the same inverse maps;
their linear parts are $-6Y$ plus constants.

The target fraction is $r^*=(d-1)/(d+1)$, already reduced for even $d$.
Its periods and directional slope are
\[
 U^*=(d+1,d-1)=\mathcal R V,\qquad
 V^*=(-2,2d)=-\mathcal R U,\qquad a^*=(d-1,-d-1).
\]
Direct substitution in (2), (9), and (15) gives
\[
\begin{array}{c|cc}
&F_i(z+U^*)-F_i(z)&F_i(z+V^*)-F_i(z)\\\hline
F_6&0&a^*\!\cdot z+1-T\\
F_8&0&a^*\!\cdot z-T.
\end{array}                                                   \tag{16}
\]
The matrix identity
$M(r^*)=\mathcal R M(1,d)\mathcal R^T+H$ proves the required growth.
The affine rotation bijects quotient classes and tile representatives,
so covering and uniqueness transfer with these specified increments.
To obtain functions on the original physical graph put
\[
 G_i(x,y)=F_i(x+1,y).                                       \tag{17}
\]
Translation by $(1,0)$ reverses parity without interchanging axes:
$\Delta_0G_i(z)=\Delta_1F_i(z+(1,0))$. It is a graph isomorphism
between these two phases, so $G_6,G_8$ are integer superharmonic and
recurrent on the original graph. Their native domains are those of
$F_i$ shifted by $-(1,0)$; growth is unchanged. Their $U^*$ constants
remain zero, while their $V^*$ constants are respectively
$1-d^2-d$ and $-d^2-d$. This is a convenient physical representative. The literal lower-left
placement instead uses $F_i(z+(0,1))$; it differs from (17) by the
even translation $(1,-1)$ and has the same recurrence and growth.
The transformed ranks exclude nonempty forbidden
sets with $y$ bounded above, and give the same strip bounds using
vertical span. No parity-changing operation has been silently treated
as a symmetry of the original graph.

\section*{7. Corrected rotated alternate statement and an exact source typo}
The following is the precise extension-and-recurrence portion of a
corrected Lemma 6.10, label \texttt{lemma:basecase\_2k\_even\_alt}.
It does not separately certify that lemma's boundary-word assertion.
For every even $d\ge4$ define
\[
\begin{split}
 A^*&=\{(-2,2d),(d+1,d-1)\},\\
 B^*&=\{(X,Y):Y-X=d,\ 0<X<d-1\},\\
 C^*&=\{(X,Y):-2\le X\le d+1,\ 1\le Y\le3d-2,
                         \ -1\le Y-X\le2d+1\},\\
 D^*_{\rm del}&=\{X\le-1,\ Y\le d-1\}
                  \cup\{X\ge d,\ Y\ge2d\},\\
 \widehat D^*&=(A^*\cup B^*\cup C^*)\setminus D^*_{\rm del}.
\end{split}                                                   \tag{18}
\]
There is a unique integer function $F_8$ with native values
\[
 -\left\lfloor\frac{(Y-X)^2}{4}\right\rfloor
 +\min(d-Y,0)+\min(2d-1-Y,0)+\max(Y-X-d,0)
 +\ind_{A^*\cup B^*}                                       \tag{19}
\]
and the specified periods (16). It is superharmonic and recurrent on
bit $1$; $G_8$ in (17) is so on the original bit $0$ graph.

To verify the domain and values, substitute
$(x,y)=(2d-1-Y,X+1)$ in (8). The lower deleted strip
$y\le0,x\ge d$ becomes $X\le-1,Y\le d-1$, and the upper one
becomes $X\ge d,Y\ge2d$. The rectangle, diagonal band, correction
line and corners become exactly (18). The elementary identities
$\min(-a,0)-\min(a,0)=-a$ and
$\max(a,0)-\max(-a,0)=a$ give (19) from (15).
This proves the corrected statement using the already established
extension and rank, with specified constants rather than an unspecified
periodicity condition.

The paper's line 2497 instead uses $Y\le d+1$ in the first deleted
strip. This removes valid points: for every even $d\ge4$, the point
$(-1,d)$ belongs to (18), and maps to the valid B7 point $(d-1,0)$.
It is not one of the exceptional corners, so the wrong deletion cannot
be repaired by re-adding a corner. The exact correction is
$Y\le d+1\longrightarrow Y\le d-1$ in that strip.
The same original statement also has an unhatted $g_{q_1}(x)$ on the
right of its alternate periodicity equation; the corrected equation
must use the same alternate function on both sides, as in (16).

The phase correction is substantive. Put $d=2m$. At the three points
$(-1,d),(0,d),(1,d)$, formula (19) has values
\[
 -m^2-m+1,\qquad -m^2,\qquad -m^2+m.
\]
Their horizontal second difference is $+1$ at the even center $(0,d)$.
Thus the unshifted formula is not superharmonic on the original bit $0$
graph, for any permitted $d$. It is correct on bit $1$, and the explicit
odd translation (17) fixes the physical phase. At the smallest parameter
$d=4$ the values are $-5,-4,-2$, making the defect concrete.

\begin{sloppypar}
The inventory identifiers of the recurrence replacements remain
Lemma 6.15, \texttt{lemma:doubled-staircase-recurrent}, and Lemma 6.16,
\texttt{lemma:doubled-alt-staircase-recurrent}; their stated fractions,
parameters and bit-$0$ formulas require no change. The auxiliary defect
belongs to Lemma 6.10, not to those two statements. None of these base
proofs establishes the final general recurrence argument's hereditary
column claims or a safe rule for using a parent rank across a child seam.
\end{sloppypar}
\end{document}
